% આ ભાગનો પૂર્ણ સંપાદનીય પાઠ છે. શૈલી, ફોન્ટ, ચિત્રો અને પ્રકરણસંદર્ભ માટે સંપૂર્ણ સ્રોત ZIP તથા reader/full-reader.tex જુઓ. આ એકલી સ્વતંત્ર કમ્પાઇલ થતી મુખ્ય ફાઇલ નથી.

% BEGIN OLP-0475
\olpart{aml}{પ્રયુક્ત મોડલ તર્ક}
\phantomsection\label{guunit:OLP-0475}


\begin{editorial}
આ ભાગમાં મોડલ તર્કના કેટલાક ઉપયોગો, જેમ કે સમયસંબંધી અને
જ્ઞાનસંબંધી તર્ક, વિશેની પ્રયોગાત્મક મુસદ્દા-સામગ્રી છે.
\end{editorial}

\begingroup
\makeatletter\def\ol@sectioncs{section}\makeatother

% BEGIN OLP-0476
\olchapter{aml}{tl}{સમયસંબંધી તર્કપદ્ધતિઓ}
\olchapterid{aml}{tl}
\phantomsection\label{guunit:OLP-0476}


\begin{editorial}
  આ પ્રકરણ સમયસંબંધી તર્કપદ્ધતિઓ વિશે છે.
\end{editorial}

\begingroup
\makeatletter\def\ol@sectioncs{section}\makeatother

% BEGIN OLP-0477
\olfileid{aml}{tl}{int}

\olsection{પરિચય}
\phantomsection\label{guunit:OLP-0477}


સમયસંબંધી તર્કપદ્ધતિઓ એવી બાબતો વિશેના દાવાઓનું અધ્યયન કરે છે જે
આગળ સાચી થશે અથવા અગાઉ સાચી હતી. સમયસંબંધી તર્કના પ્રણેતા તરીકે
આર્થર પ્રાયરને શ્રેય આપવામાં આવે છે; તેમણે તેને \emph{કાળ-તર્ક}
કહ્યો હતો. અહીં અમે સમયસંબંધી તર્કનો અભ્યાસ મોટાભાગે પ્રાયરના મૂળ
મોડલ અભિગમ મુજબ કરીશું: તેમાં વિધાનતર્કના મૂળભૂત માળખામાં
સમયસંબંધી સંક્રિયકો ઉમેરાય છે, જ્યારે વિધાનતર્ક સામાન્ય રીતે
વિધાનોને કાળ-નિર્દેશ વિનાનાં ગણે છે.

દાખલા તરીકે, વિધાનતર્કમાં હું બીઝી નામના કૂતરા વિશે વાત કરી શકું:
કૂતરાંઓની ટેવ મુજબ તે ક્યારેક બેસે છે અને ક્યારેક બેસતો નથી.
શાસ્ત્રીય તર્કમાં બીઝી બેસે છે અને બીઝી બેસતો નથી, એવા બંને દાવા
એકસાથે કરવા વિરોધાભાસી ગણાય. પરંતુ બંને દાવા સાચા હોઈ શકે છે,
માત્ર એક જ સમયે નહીં; ભાષામાં સમયસંબંધી સંક્રિયકો ઉમેરવાથી આ વાત
સરળતાથી વ્યક્ત કરી શકાય છે. આ સંક્રિયકો ઉમેરવાથી ``બીઝીને આગળ
ખાવાનું ઇનામ અથવા દડો મળશે'' પરથી ``બીઝીને આગળ ખાવાનું ઇનામ
મળશે અથવા બીઝીને આગળ દડો મળશે'' જેવા અનુમાનોની માન્યતા પણ
સમજાવી શકાય છે.

જોકે સમયસંબંધી તર્કમાં ઘણા દાર્શનિક પ્રશ્નો ઊભા થાય છે, જેને કારણે
આપણે તેની એક રૂપરેખા બીજી કરતાં પસંદ કરી શકીએ. દાખલા તરીકે,
ભવિષ્યસંબંધી આકસ્મિક વિધાન ભવિષ્ય વિશેનું એવું વિધાન છે જે ન તો
અનિવાર્ય છે ન તો અશક્ય. જો આપણે કહીએ કે ``રિચાર્ડ આવતી કાલે
કરિયાણાની દુકાને જશે'', તો આપણે હજી ન બનેલી બાબત વિશે દાવો કરીએ
છીએ, જેનું સત્યમૂલ્ય વિવાદાસ્પદ છે. હકીકતમાં, એ દાવાને મૂળે
સત્યમૂલ્ય \emph{આપી} પણ શકાય છે કે કેમ એ વિવાદાસ્પદ છે. જો આપણે
કડક નિયતિવાદી હોઈએ, તો સંબંધિત ઘટના બનવાની હોય તે પહેલાં પણ આ
વાક્ય હકીકતમાં સાચું કે ખોટું છે, એવું માનવામાં કદાચ આપણને વાંધો
ન હોય---માત્ર તેનું સત્યમૂલ્ય હજી આપણે જાણતા ન હોઈએ. તેનાથી
ઊલટું, આપણે એવું માની શકીએ કે ભવિષ્ય ખરેખર ખુલ્લું છે અને
ભવિષ્યસંબંધી આકસ્મિક વિધાનોના સત્યમૂલ્યો અનિર્ધારિત છે.

સમયની રચના અને સ્વરૂપ અંગેની આવી ઘણી માન્યતાઓ સમયસંબંધી
તર્કપદ્ધતિઓનાં નિદર્શનો અને રૂપરેખાઓની આપણી પસંદગીમાં વણાયેલી હોય છે.
દાખલા તરીકે, આપણે પૂછવું પડે કે સમય રેખીય, શાખાવાળો કે ચક્રીય
હોય એવા નિદર્શનો રચવા જોઈએ? આપણાં સમયસંબંધી નિદર્શનોમાં આરંભબિંદુ અને
અંતબિંદુ હોવા જોઈએ કે નહીં, અને સમયને વિવિક્ત ક્ષણો દ્વારા
દર્શાવવો કે સાતત્ય તરીકે, એ પણ નક્કી કરવું પડે.
% END OLP-0477
\endgroup

\begingroup
\makeatletter\def\ol@sectioncs{section}\makeatother

% BEGIN OLP-0478
\olfileid{aml}{tl}{sem}

\olsection{સમયસંબંધી તર્કનું અર્થવિજ્ઞાન}
\phantomsection\label{guunit:OLP-0478}


\begin{defn}
સમયસંબંધી તર્કની મૂળભૂત ભાષામાં નીચેની બાબતો હોય છે:
\begin{enumerate}
  \item અસત્ય માટેનો વિધાનાત્મક અચલાંક~$\lfalse$.
  \item સત્ય માટેનો વિધાનાત્મક અચલાંક~$\ltrue$.
  \item પ્રતિજ્ઞાપન ચલોનો ગણનીય ગણ: $\Obj
    p_0$, $\Obj p_1$, $\Obj p_2$, \dots
  \item વિધાનાત્મક સંયોજકો: \startycommalist
  \ycomma $\lnot$ (નિષેધ)%
  \ycomma $\land$ (સંયોજન)%
  \ycomma $\lor$ (વિયોજન)%
  \ycomma $\lif$ (પ્રેરણ)%
  \ycomma $\liff$ (દ્વિમુખી પ્રેરણ).
  \item ભૂતકાળના સંક્રિયકો $\Ptemp$ અને $\Htemp$.
  \item ભવિષ્યના સંક્રિયકો $\Ftemp$ અને $\Gtemp$.
\end{enumerate}
\end{defn}

આગળ જઈને બીજા પ્રકારના મોડલ સંક્રિયકો ઉમેરવાની શક્યતા વિશે ચર્ચા કરીશું.

\begin{defn}
સમયસંબંધી ભાષાનાં \emph{સૂત્રો} નીચે મુજબ આગમનાત્મક રીતે
  વ્યાખ્યાયિત થાય છે:
\begin{enumerate}
\item $\lfalse$ પરમાણુક સૂત્ર છે.

\item $\ltrue$ પરમાણુક સૂત્ર છે.

\item દરેક વિધાનાત્મક ચલ $\Obj p_i$ (પરમાણુક) સૂત્ર છે.

\item જો $!A$ એ સૂત્ર હોય, તો $\lnot !A$ પણ
  સૂત્ર છે.

\item જો $!A$ અને $!B$ એ સૂત્રો હોય, તો $(!A \land
  !B)$ પણ સૂત્ર છે.

\item જો $!A$ અને $!B$ એ સૂત્રો હોય, તો $(!A \lor !B)$
  પણ સૂત્ર છે.

\item જો $!A$ અને $!B$ એ સૂત્રો હોય, તો $(!A \lif !B)$
  પણ સૂત્ર છે.

\item જો $!A$ અને $!B$ એ સૂત્રો હોય, તો $(!A \liff !B)$
  પણ સૂત્ર છે.

\sourcecorrection{OLMD-076}{મૂળમાં આ યાદીના ભવિષ્ય-સંક્રિયકને
  $F$ લખ્યો છે; ભાષાની અગાઉની યાદી અને આગળની અર્થવ્યાખ્યા મુજબ
  તેને $\Ftemp$ કર્યો છે.}
\item જો $!A$ એ સૂત્ર હોય, તો $\Ptemp  !A$, $\Htemp  !A$, $\Ftemp !A$, $\Gtemp  !A$
  એ બધાં સૂત્રો છે.

\item આ સિવાય બીજું કશું સૂત્ર નથી.
\end{enumerate}
\end{defn}

બીજી ઇન્ટેન્શનલ તર્કપદ્ધતિઓની જેમ સમયસંબંધી તર્કપદ્ધતિઓનું
અર્થવિજ્ઞાન પણ સંબંધાત્મક નિદર્શનો વડે આપવામાં આવે છે.

\begin{defn}
  સમયસંબંધી ભાષાનું \emph{નિદર્શ} એ ત્રિક
  $\mModel{M} = \tuple{T, \prec, V}$ છે, જ્યાં
  \begin{enumerate}
  \item $T$ સમયબિંદુઓ તરીકે અર્થઘટિત થતો અખાલી ગણ છે.
  \item $\prec$ એ~$T$ પરનો દ્વિસ્થાની સંબંધ છે.
  \item $V$ એવું વિધેય છે જે દરેક પ્રતિજ્ઞાપન ચલ~$p$ને સમયબિંદુઓનો ગણ $V(p)$ આપે છે.
  \end{enumerate}
  $t \prec t'$ હોય ત્યારે કહીએ કે $t$ એ~$t'$ની \emph{પહેલાં આવે છે}.
  $t \in V(p)$ હોય ત્યારે કહીએ કે $p$ એ~$t$એ \emph{સત્ય છે}.
\end{defn}

હમણાં આપણે પૂર્વવર્તિતાના સંબંધ~$\prec$ પર કોઈ શરત મૂકતા નથી.
એટલે હાલ આપણાં સમયસંબંધી નિદર્શનોની રચના પર કોઈ નિયંત્રણ નથી:
સમય રેખીય, શાખાવાળો, ચક્રીય કે અન્ય કોઈ પણ રચનાવાળો હોય એવા
નિદર્શનો હોઈ શકે છે.

સામાન્ય મોડલ તર્કની જેમ દરેક સમયસંબંધી નિદર્શ એ નક્કી કરે છે કે
તેના કયા બિંદુએ કયાં સૂત્રો સત્ય ગણાય. સંબંધાત્મક
અર્થવિજ્ઞાનના મૂળ ખ્યાલ માટે આપણે એ જ સંજ્ઞા વાપરીએ છીએ:
``નિદર્શ~$\mModel{M}$માં સૂત્ર~$!A$ બિંદુ~$t$એ સત્ય છે.''
આ સંબંધ આગમનાત્મક રીતે વ્યાખ્યાયિત થાય છે અને બધા બિનમોડલ
સંક્રિયકો માટે સામાન્ય મોડલ તર્કના કિસ્સા જેવો જ છે.

\begin{defn}\ollabel{defn:tmodels}
  નિદર્શ~$\mModel M$માં \emph{સૂત્ર~$!A$નું~$t$એ સત્ય હોવું},
  સંજ્ઞામાં $\mSat{M}{!A}[t]$, નીચે મુજબ આગમનાત્મક રીતે
  વ્યાખ્યાયિત થાય છે:
  \begin{enumerate}
  \item \indcase{!A}{\lfalse}{$\mSat{M}{\lfalse}[t]$ ક્યારેય નહીં}.
  \item \indcase{!A}{\ltrue}{$\mSat{M}{\ltrue}[t]$ હંમેશા}.
  \item $\mSat{M}{p}[t]$ ત્યારે અને માત્ર ત્યારે જ જ્યારે $t \in V(p)$
  \item \indcase{!A}{\lnot !B}{$\mSat{M}{\indfrm}[t]$ ત્યારે અને માત્ર ત્યારે જ જ્યારે
    $\mSat/{M}{!B}[t]$}.
  \item \indcase{!A}{(!B \land !C)}{$\mSat{M}{\indfrm}[t]$ ત્યારે અને માત્ર ત્યારે જ જ્યારે
    $\mSat{M}{!B}[t]$ અને $\mSat{M}{!C}[t]$}.
  \item \indcase{!A}{(!B \lor !C)}{$\mSat{M}{\indfrm}[t]$ ત્યારે અને માત્ર ત્યારે જ જ્યારે
    $\mSat{M}{!B}[t]$ અથવા $\mSat{M}{!C}[t]$} (અથવા બંને).
  \item \indcase{!A}{(!B \lif !C)}{$\mSat{M}{\indfrm}[t]$ ત્યારે અને માત્ર ત્યારે જ જ્યારે
    $\mSat/{M}{!B}[t]$ અથવા $\mSat{M}{!C}[t]$}.
  \item \indcase{!A}{(!B \liff !C)}{$\mSat{M}{\indfrm}[t]$ ત્યારે અને માત્ર ત્યારે જ જ્યારે
    $\mSat{M}{!B}[t]$ અને $\mSat{M}{!C}[t]$ બંને હોય, અથવા
    $\mSat{M}{!B}[t]$ તથા $\mSat{M}{!C}[t]$ બેમાંથી એકેય ન હોય}.
  \item\ollabel{defn:sub:mmodels-p}
    \indcase{!A}{\Ptemp  !B}{$\mSat{M}{\indfrm}[t]$ ત્યારે અને માત્ર ત્યારે જ જ્યારે
    $t' \prec t$ ધરાવતા કોઈ $t' \in T$ માટે $\mSat{M}{!B}[t']$}
\item\ollabel{defn:sub:mmodels-h}
    \indcase{!A}{\Htemp  !B}{$\mSat{M}{\indfrm}[t]$ ત્યારે અને માત્ર ત્યારે જ જ્યારે
    $t' \prec t$ ધરાવતા દરેક $t' \in T$ માટે $\mSat{M}{!B}[t']$}
   \item\ollabel{defn:sub:mmodels-f}
    \indcase{!A}{\Ftemp  !B}{$\mSat{M}{\indfrm}[t]$ ત્યારે અને માત્ર ત્યારે જ જ્યારે
    $t \prec t'$ ધરાવતા કોઈ $t' \in T$ માટે $\mSat{M}{!B}[t']$}
\item\ollabel{defn:sub:mmodels-g}
    \indcase{!A}{\Gtemp  !B}{$\mSat{M}{\indfrm}[t]$ ત્યારે અને માત્ર ત્યારે જ જ્યારે
    $t \prec t'$ ધરાવતા દરેક $t' \in T$ માટે $\mSat{M}{!B}[t']$}
  \end{enumerate} 
\end{defn}

આ અર્થવિજ્ઞાન પરથી જોઈ શકાય છે કે સંક્રિયકો~$\Ptemp$ અને~$\Htemp$
પરસ્પર દ્વૈત છે, અને સંક્રિયકો~$\Ftemp$ તથા~$\Gtemp$ પણ એમ જ છે.
એટલે $\Htemp  !A$ને $\lnot \Ptemp \lnot !A$ તરીકે વ્યાખ્યાયિત
કરી શકીએ છીએ; $\Gtemp$ અને~$\Ftemp$ માટે પણ આવું જ થાય છે.
% END OLP-0478
\endgroup

\begingroup
\makeatletter\def\ol@sectioncs{section}\makeatother

% BEGIN OLP-0479
\olfileid{aml}{tl}{acc}

\olsection{સમયસંબંધી ફ્રેમોના ગુણધર્મો}
\phantomsection\label{guunit:OLP-0479}


આપણાં સમયસંબંધી નિદર્શનોમાં સંબંધ~$\prec$ પર કોઈ શરત નથી,
એટલે આપણા પરિચિત સ્વયંસિદ્ધાંતોમાંથી બધાં નિદર્શનોમાં માન્ય
રહેતો એકમાત્ર સ્વયંસિદ્ધાંત~$K$, અથવા તેના અનુરૂપ
$K_{\Gtemp}$ અને~$K_{\Htemp}$ છે:
\begin{align*}
\tag{$K_{\Gtemp}$}  \Gtemp (p \to q) & \to (\Gtemp p \to \Gtemp q)\\
\tag{$K_{\Htemp}$}  \Htemp (p \to q) & \to (\Htemp p \to \Htemp q)
\end{align*}

પરંતુ સમયને કેવી રીતે દર્શાવવો તે અંગે નિદર્શનો પર વધુ કડક
શરતો મૂકીએ, દાખલા તરીકે~$\prec$ રેખીય ક્રમ હોય તે સુનિશ્ચિત
કરીએ, તો નિદર્શનોમાં બીજા દાવાઓ પણ માન્ય થશે.

\begin{table}[t]
  \begin{tabular}{| p{.47\linewidth} || p{.47\linewidth} |}
    \hline
    {\emph{જો $\prec$ \dots હોય}} & {\emph{તો~$\mModel{M}$માં \dots સત્ય છે:}} \\
    \hline \hline
    \emph{પરંપરિત}: \newline
    $\forall u \forall v \forall w ((u \prec v \land v \prec w) \lif u \prec w)$ & 
    $\Ftemp \Ftemp p \lif \Ftemp p$  \\
    \hline 
    \emph{રેખીય}: \newline
    $\forall w \forall v (w \prec v \lor w = v \lor v \prec w)$ &  
    $(\Ftemp \Ptemp  p \lor \Ptemp \Ftemp p) \lif (\Ptemp p \lor p \lor \Ftemp p)$\\
    \hline
    \emph{સઘન}: \newline
    $\forall w \forall v (w \prec v \to \exists u(w \prec u \land u \prec v))$ &  
    $\Ftemp p \lif \Ftemp \Ftemp p$ \\
    \hline
    \emph{ભૂતકાળ તરફ અસીમ}: \newline
    $\forall w \exists v( v \prec w)$ &  
    $\Htemp p \to \Ptemp p$ \\
    \hline
    \emph{ભવિષ્ય તરફ અસીમ}: \newline
    $\forall w \exists v( w \prec v)$ &  
    $\Gtemp p \to \Ftemp p$ \\
    \hline
  \end{tabular}
  \caption{સમયસંબંધી ફ્રેમોના કેટલાક અનુરૂપતા-ગુણધર્મો.}
  \ollabel{tab:correspondence}
\end{table} 

\olref{tab:correspondence}ના કેટલાક ગુણધર્મો સમય દર્શાવવા માટેના
નિદર્શનમાં ઇચ્છનીય લાગે. જોકે સંબંધ~$\prec$ પર આપણે ગમે તેવી
શરતો મૂકી શકીએ, તેમ છતાં બધી શરતોને સમયસંબંધી તર્કની ભાષામાં
વ્યક્ત થઈ શકતાં સૂત્રો અનુરૂપ હોતાં નથી. દાખલા તરીકે,
અસ્વવાચકતા, એટલે કે $\forall w \lnot (w \prec w)$ હોવાનો
ગુણધર્મ, સમયસંબંધી તર્કમાં કોઈ અનુરૂપ સૂત્ર ધરાવતો નથી.
% END OLP-0479
\endgroup

\begingroup
\makeatletter\def\ol@sectioncs{section}\makeatother

% BEGIN OLP-0480
\olfileid{aml}{tl}{ext}

\olsection{સમયસંબંધી તર્ક માટે વધારાના સંક્રિયકો}
\phantomsection\label{guunit:OLP-0480}


ભૂતકાળ અને ભવિષ્ય માટેના એકસ્થાની સંક્રિયકો ઉપરાંત સમયસંબંધી
તર્કપદ્ધતિઓમાં ક્યારેક ``ત્યારથી'' અને ``ત્યાં સુધી'' દર્શાવતા
દ્વિસ્થાની સંક્રિયકો~$\Since $ અને~$\Until$ પણ હોય છે. એટલે
સમયસંબંધી તર્કની ભાષામાં $\Since $ અને~$\Until$ ઉમેરીએ છીએ અને
સમયસંબંધી સૂત્રની વ્યાખ્યામાં નીચેની શરત ઉમેરીએ છીએ:

\begin{itemize}
\item[] જો $!A$ અને $!B$ એ સૂત્રો હોય, તો $(\Since  !A !B)$ અને
  $(\Until !A !B)$ બંને સૂત્રો છે.
\end{itemize}

પછી આ સંક્રિયકોનું અર્થવિજ્ઞાન નીચે મુજબ છે:

\begin{defn}\ollabel{defn:since-until}
  નિદર્શ~$\mModel M$માં \emph{સૂત્ર~$!A$નું~$t$એ સત્ય હોવું}:
  \begin{enumerate}
  \item\ollabel{defn:sub:mmodels-since} \indcase{!A}{\Since  !B
    !C}{$\mSat{M}{\indfrm}[t]$ ત્યારે અને માત્ર ત્યારે જ જ્યારે
    $t' \prec t$ ધરાવતા કોઈ $t' \in T$ માટે $\mSat{M}{!B}[t']$ હોય,
    અને $t' \prec s \prec t$ ધરાવતા દરેક $s$ માટે
    $\mSat{M}{!C}[s]$ હોય}
   \item\ollabel{defn:sub:mmodels-until} \indcase{!A}{\Until !B
    !C}{$\mSat{M}{\indfrm}[t]$ ત્યારે અને માત્ર ત્યારે જ જ્યારે
    $t \prec t'$ ધરાવતા કોઈ $t' \in T$ માટે $\mSat{M}{!B}[t']$ હોય,
    અને $t \prec s \prec t'$ ધરાવતા દરેક $s$ માટે
    $\mSat{M}{!C}[s]$ હોય}
  \end{enumerate} 
\end{defn}

$\Since  !B !C$નો સહજ અર્થ છે: ``$!B$ સાચું હતું ત્યારથી
$!C$~સાચું રહ્યું છે.'' અને $\Until !B !C$નો સહજ અર્થ છે:
``$!B$ સાચું થશે ત્યાં સુધી $!C$~સાચું રહેશે.''
% END OLP-0480
\endgroup

\begingroup
\makeatletter\def\ol@sectioncs{section}\makeatother

% BEGIN OLP-0481
\olfileid{aml}{tl}{poss}

\olsection{સંભવિત ઇતિહાસો}
\phantomsection\label{guunit:OLP-0481}


અત્યાર સુધી વાપરેલાં સમયસંબંધી તર્કનાં સંબંધાત્મક નિદર્શનો ઘણાં
લવચીક છે, કારણ કે સુલભતા સંબંધ પર કોઈ નિયંત્રણ મૂકવું પડતું નથી.
એટલે સમયસંબંધી નિદર્શનોમાં ભૂતકાળ અને ભવિષ્ય બંને તરફ શાખાઓ પડી
શકે છે. પરંતુ આપણે શાખાઓનો વધુ ``મોડલ'' ખ્યાલ લઈ શકીએ, જેમાં
ઘટનાઓની શ્રેણીઓને સંભવિત ઇતિહાસો ગણીએ. આ માટે ભાષા બદલવી જરૂરી
નથી; જોકે આપણે આપણા ``સામાન્ય'' મોડલ સંક્રિયકો~$\Box$ અને~$\Diamond$
પણ ઉમેરી શકીએ, તથા સમય સાથે કર્તાઓના જ્ઞાનમાં થતા ફેરફારો દર્શાવવા
જ્ઞાનસંબંધી સુલભતા સંબંધો ઉમેરવાનું પણ વિચારી શકીએ.

\begin{defn}
  સમયસંબંધી ભાષાનું \emph{સંભવિત ઇતિહાસોનું નિદર્શ} એ ત્રિક
  $\mModel{M} = \tuple{T, C, V}$ છે, જ્યાં
  \begin{enumerate}
    \item $T$ સમયની સ્થિતિઓ તરીકે અર્થઘટિત થતો અખાલી ગણ છે.
    \item $C$ એ તંત્રના સંગણનાત્મક પથોનો, એટલે કે \emph{સંભવિત
      ઇતિહાસો}નો ગણ છે. બીજા શબ્દોમાં, $C$ એ સ્થિતિઓ
      $s_1$, $s_2$,~$s_3$, \dots ની શ્રેણીઓ~$\sigma$નો ગણ છે,
      જ્યાં દરેક $s_i \in T$ છે.
    \item $V$ એવું વિધેય છે જે દરેક પ્રતિજ્ઞાપન ચલ~$p$ને સમયબિંદુઓનો ગણ~$V(p)$ આપે છે.
  \end{enumerate}
  સરળતા માટે સામાન્ય રીતે એવું પણ માનીએ છીએ કે કોઈ ઇતિહાસ~$C$માં
  હોય, તો તેની બધી અંતિમ ઉપશ્રેણીઓ પણ તેમાં હોય. દાખલા તરીકે,
  $s_1$, $s_2$,~$s_3$ શ્રેણી~$C$માં હોય, તો $s_2$, $s_3$ અને~$s_3$
  પણ તેમાં હોય. વધુમાં, બે સ્થિતિઓ $s_i$ અને~$s_j$ શ્રેણી~$\sigma$માં
  આવે ત્યારે, $i < j$ હોય તો $s_i \prec_\sigma s_j$ કહીએ છીએ.
  $t \in V(p)$ હોય ત્યારે કહીએ કે $p$ એ~$t$એ \emph{સત્ય છે}.
\end{defn}

અહીં સંબંધિત એક ફેરફાર એ છે કે નિદર્શ~$\mModel M$માં
સમયબિંદુ~$t$એ સૂત્રનું સત્ય ચકાસીએ ત્યારે તેને એવા
ઇતિહાસ~$\sigma$ને અનુલક્ષીને ચકાસીએ છીએ જેમાં~$t$ એક સ્થિતિ તરીકે
આવે છે. જોકે પ્રતિજ્ઞાપન ચલો અથવા સત્યવિધેયી
સંયોજકોનું અર્થવિજ્ઞાન બદલવાની જરૂર નથી. તેઓ~$\sigma$નો ઉલ્લેખ
કરતા નથી, એટલે તેમના નિયમો \olref[sem]{defn:tmodels} જેવા જ છે.
પરંતુ હવે આ ઇતિહાસોને અનુલક્ષીને ભવિષ્ય-સંક્રિયક~$\Ftemp$ને ફરી
વ્યાખ્યાયિત કરીએ છીએ અને સંક્રિયક~$\Diamond$ ઉમેરીએ છીએ.

\begin{defn}\ollabel{defn:phmodels}
  નિદર્શ~$\mModel M$માં \emph{સૂત્ર~$!A$નું~$t, \sigma$એ સત્ય હોવું},
  સંજ્ઞામાં $\mSat{M}{!A}[t, \sigma]$:
  \begin{enumerate}
  \item\ollabel{defn:sub:phmodels-f} \indcase{!A}{\Ftemp
    !B}{$\mSat{M}{\indfrm}[t, \sigma]$ ત્યારે અને માત્ર ત્યારે જ જ્યારે
    $t \prec_\sigma t'$ ધરાવતા કોઈ $t' \in T$ માટે
    $\mSat{M}{!B}[t', \sigma]$.}
  \item\ollabel{defn:sub:phmodels-diamond} \indcase{!A}{\Diamond
    !B}{$\mSat{M}{\indfrm}[t, \sigma]$ ત્યારે અને માત્ર ત્યારે જ જ્યારે
    $t$ આવતું હોય એવા કોઈ $\sigma' \in C$ માટે
    $\mSat{M}{!B}[t, \sigma']$.}
  \end{enumerate} 
\end{defn}

બીજા સમયસંબંધી અને મોડલ સંક્રિયકો પણ એ જ રીતે વ્યાખ્યાયિત કરી
શકાય. પરંતુ હવે કાળ અને શક્યતાને જોડતા દાવાઓ વ્યક્ત કરી શકીએ છીએ.
દાખલા તરીકે, ``$p$~બનશે નહીં, પણ તે બની શક્યું હોત'' એ વાત
$\lnot \Ftemp p \land \Diamond \Ftemp p$ સૂત્રથી દર્શાવી શકીએ.
આ સૂત્ર એવા બિંદુ અને ઇતિહાસમાં સત્ય હશે જ્યાં પછીની કોઈ પણ સ્થિતિએ
$p$ સત્ય થતું નથી, પરંતુ કોઈ વૈકલ્પિક ઇતિહાસમાં $p$ આગળ સત્ય થશે.
% END OLP-0481
\endgroup


\OLEndPartHook
% END OLP-0476
\endgroup


\begingroup
\makeatletter\def\ol@sectioncs{section}\makeatother

% BEGIN OLP-0482
\olchapter{aml}{el}{જ્ઞાનસંબંધી તર્કપદ્ધતિઓ}
\olchapterid{aml}{el}
\phantomsection\label{guunit:OLP-0482}


\begin{editorial}
  આ પ્રકરણ જ્ઞાનસંબંધી તર્કપદ્ધતિઓના અધિસિદ્ધાંત વિશે છે. તેની
  રચના મૂળભૂત શાસ્ત્રીય મોડલ તર્ક અંગેની આલ્ડો એન્ટોનેલીની નોંધો
  જેવી છે, પરંતુ બિસિમ્યુલેશન અને ગતિશીલ જ્ઞાનસંબંધી તર્કપદ્ધતિઓ
  અંગે સામગ્રી ઉમેરવા ઓડ્રી યેપે તેને ફરી લખ્યું છે.
\end{editorial}

\begingroup
\makeatletter\def\ol@sectioncs{section}\makeatother

% BEGIN OLP-0483
\olfileid{aml}{el}{int}

\olsection{પરિચય}
\phantomsection\label{guunit:OLP-0483}


મોડલ તર્ક \emph{મોડલ વિધાનો} અને તેમની વચ્ચેના તાર્કિક પરિણામોના
સંબંધોનો અભ્યાસ કરે છે; તે જ રીતે જ્ઞાનસંબંધી તર્ક
\emph{જ્ઞાનસંબંધી વિધાનો} અને તેમની વચ્ચેના તાર્કિક પરિણામોના
સંબંધોનો અભ્યાસ કરે છે. મોડલ સંક્રિયકોને શક્યતા અને અનિવાર્યતા
દર્શાવતા ગણવાને બદલે, અહીં એકસ્થાની સંયોજકોનું અર્થઘટન જ્ઞાન અને
માન્યતાના નિદર્શન માટે જ્ઞાનસંબંધી અથવા માન્યતાસંબંધી રીતે થાય છે.
દાખલા તરીકે, આપણે નીચેના જેવા દાવા વ્યક્ત કરવા માગી શકીએ:

\begin{enumerate}
\item રિચાર્ડ જાણે છે કે કૅલ્ગરી આલ્બર્ટામાં છે.
\item ઓડ્રીને લાગે છે કે સોફા પર કૂતરો હોઈ શકે છે.
\item રિચાર્ડ જાણે છે કે ઓડ્રી જાણે છે કે તેનો વર્ગ મંગળવારે હોય છે.
\item દરેક જણ જાણે છે કે વર્ષમાં બાર મહિના હોય છે.
\end{enumerate}
\sourcecorrection{OLMD-077}{મૂળમાં લેખકનું નામ ``Jaako Hintikka''
લખ્યું છે; ગ્રંથસૂચિ-વિવરણ મુજબ ``Jaakko Hintikka'' છે. અહીં નામ
સુધાર્યું છે.}
આધુનિક જ્ઞાનસંબંધી તર્કની શરૂઆત ઘણી વાર યાક્કો હિન્તિક્કાના
૧૯૬૨ના પુસ્તક~\emph{Knowledge and Belief}માં શોધાય છે. એ પુસ્તક
લખાયું ત્યારે તર્કમાં શક્ય જગતોનું અર્થવિજ્ઞાન વધુને વધુ વપરાવા
લાગ્યું હતું. હકીકતમાં, જ્ઞાનસંબંધી તર્કપદ્ધતિઓ બીજી મોડલ
તર્કપદ્ધતિઓનાં મોટા ભાગનાં એ જ અર્થવિજ્ઞાનલક્ષી સાધનો વાપરે છે,
પરંતુ તેમનું અર્થઘટન જુદી રીતે કરે છે. મુખ્ય ફેરફાર એ છે કે
\emph{સુલભતા સંબંધ} શું દર્શાવે છે એવું આપણે માનીએ છીએ.
જ્ઞાનસંબંધી તર્કમાં તે કોઈ પ્રકારની \emph{જ્ઞાનસંબંધી શક્યતા}
દર્શાવે છે. આપણે જોઈશું કે જે જ્ઞાનસંબંધી ખ્યાલનું નિદર્શન કરીએ
તેના પર સુલભતા સંબંધને મૂકવાની શરતો આધારિત હોય છે. અનુરૂપતા
સિદ્ધાંતને જ્ઞાનસંબંધી અર્થઘટન આપીએ ત્યારે શું થાય છે તે પણ
જોઈશું. ઉપરનાં ઉદાહરણોમાં બે કર્તાઓ, રિચાર્ડ અને ઓડ્રી, તથા
દરેક શું જાણે છે તેનો પરસ્પર સંબંધ આવે છે. આપણે વિચારીશું તે
જ્ઞાનસંબંધી તર્કપદ્ધતિઓ બહુ-કર્તા તર્કપદ્ધતિઓ છે, જેમાં આવી બાબતો
વ્યક્ત થઈ શકે છે. તેની સામે, એક-કર્તા જ્ઞાનસંબંધી તર્કપદ્ધતિ
માત્ર એક વ્યક્તિ શું જાણે છે કે માને છે તે વિશે વાત કરે છે.
% END OLP-0483
\endgroup

\begingroup
\makeatletter\def\ol@sectioncs{section}\makeatother

% BEGIN OLP-0484
\olfileid{aml}{el}{lan}

\olsection{જ્ઞાનસંબંધી તર્કની ભાષા}
\phantomsection\label{guunit:OLP-0484}


\begin{defn}
$G$ એ કર્તા-સંકેતોનો ગણ હોય. બહુ-કર્તા જ્ઞાનસંબંધી તર્કની
મૂળભૂત ભાષામાં નીચેની બાબતો હોય છે:
\begin{enumerate}
  \item અસત્ય માટેનો વિધાનાત્મક અચલાંક~$\lfalse$.
  \item સત્ય માટેનો વિધાનાત્મક અચલાંક~$\ltrue$.
  \item પ્રતિજ્ઞાપન ચલોનો ગણનીય ગણ: $\Obj
    p_0$, $\Obj p_1$, $\Obj p_2$, \dots
  \item વિધાનાત્મક સંયોજકો: \startycommalist
  \ycomma $\lnot$ (નિષેધ)%
  \ycomma $\land$ (સંયોજન)%
  \ycomma $\lor$ (વિયોજન)%
  \ycomma $\lif$ (પ્રેરણ)%
  \ycomma $\liff$ (દ્વિમુખી પ્રેરણ).
  \item $a \in G$ હોય ત્યારે જ્ઞાન-સંક્રિયક $\Knows_a$.
\end{enumerate}
\end{defn}

જો આપણાં તંત્રમાં માત્ર એક જ કર્તાના જ્ઞાનનો વિચાર કરીએ, તો
ગણ~$G$ અને જુદા જુદા કર્તાઓનો ઉલ્લેખ છોડી શકીએ. એ કિસ્સામાં
માત્ર મૂળભૂત સંક્રિયક~$\Knows$ રહે છે.

\begin{defn}
જ્ઞાનસંબંધી ભાષાનાં \emph{સૂત્રો} નીચે મુજબ આગમનાત્મક રીતે
  વ્યાખ્યાયિત થાય છે:
\begin{enumerate}
\item $\lfalse$ પરમાણુક સૂત્ર છે.

\item $\ltrue$ પરમાણુક સૂત્ર છે.

\item દરેક વિધાનાત્મક ચલ $\Obj p_i$ (પરમાણુક) સૂત્ર છે.

\item જો $!A$ એ સૂત્ર હોય, તો $\lnot !A$ પણ
  સૂત્ર છે.

\item જો $!A$ અને $!B$ એ સૂત્રો હોય, તો $(!A \land
  !B)$ પણ સૂત્ર છે.

\item જો $!A$ અને $!B$ એ સૂત્રો હોય, તો $(!A \lor !B)$
  પણ સૂત્ર છે.

\item જો $!A$ અને $!B$ એ સૂત્રો હોય, તો $(!A \lif !B)$
  પણ સૂત્ર છે.

\item જો $!A$ અને $!B$ એ સૂત્રો હોય, તો $(!A \liff !B)$
  પણ સૂત્ર છે.

\item જો $!A$ એ સૂત્ર હોય અને $a \in G$ હોય, તો $\Knows_a !A$
  પણ સૂત્ર છે.

\item આ સિવાય બીજું કશું સૂત્ર નથી.
\end{enumerate}
\end{defn}

જો સૂત્ર~$!A$માં $\Knows_a$ ન હોય, તો તેને
\emph{મોડલ-મુક્ત} કહીએ છીએ.

\begin{defn}
  સંક્રિયક~$\Knows$ વ્યક્તિગત જ્ઞાન દર્શાવવા માટે છે, જ્યારે
  $\EKnows$, જેને ઘણી વાર ``દરેક જણ જાણે છે'' એમ વાંચીએ છીએ,
  સમૂહ-જ્ઞાન દર્શાવે છે. $G' \subseteq G$ હોય ત્યારે
  $\EKnows_{G'} !A$ને \[\bigwedge_{b \in G'} \Knows_b !A.\]
  માટેના સંક્ષેપ તરીકે વ્યાખ્યાયિત કરીએ છીએ.
\end{defn}

જ્ઞાનનો એથી પણ વધુ પ્રબળ ખ્યાલ, એટલે કે કર્તાઓના ગણ~$G$માં
\emph{સામાન્ય જ્ઞાન}, પણ વ્યાખ્યાયિત કરી શકીએ છીએ. કોઈ માહિતી
કર્તાઓના સમૂહમાં સામાન્ય જ્ઞાન હોય, તો તેનો અર્થ એ છે કે એ સમૂહના
કર્તાઓના દરેક સંયોજન માટે સૌ જાણે છે કે બીજાં જાણે છે કે બીજાં
જાણે છે \dots એમ અનંત સુધી. આ સમૂહ-જ્ઞાન કરતાં ઘણું વધુ પ્રબળ
છે; એવા સંબંધાત્મક નિદર્શનો સરળતાથી રચી શકાય છે જેમાં
સૂત્ર સમૂહ-જ્ઞાન હોય પરંતુ સામાન્ય જ્ઞાન ન હોય.
``$G$માં~$!A$ સામાન્ય જ્ઞાન છે'' એ દર્શાવવા
$\CKnows_G !A$ વાપરીશું.
% END OLP-0484
\endgroup

\begingroup
\makeatletter\def\ol@sectioncs{section}\makeatother

% BEGIN OLP-0485
\olfileid{aml}{el}{rel}

\olsection{સંબંધાત્મક નિદર્શનો}
\phantomsection\label{guunit:OLP-0485}


જ્ઞાનસંબંધી તર્કપદ્ધતિઓનો મૂળભૂત અર્થવિજ્ઞાનલક્ષી ખ્યાલ સામાન્ય
મોડલ તર્કપદ્ધતિઓ જેવો જ છે. સંબંધાત્મક નિદર્શમાં હજી પણ જગતોનો
ગણ અને કયા જગતમાં કયાં પ્રતિજ્ઞાપન ચલોને ``સત્ય'' ગણવા તે નક્કી કરતી નિમણૂક હોય છે. એક જ
કર્તાનો વિચાર કરતા હોઈએ, તો હંમેશની જેમ એક જ સુલભતા સંબંધ
હોય છે. પરંતુ બહુ-કર્તા જ્ઞાનસંબંધી તર્કમાં એ સંબંધની જગ્યાએ
સુલભતા સંબંધોનો ગણ આવે છે: કર્તા-સંકેતોના ગણ~$G$માંના દરેક~$a$
માટે એક સંબંધ.

\emph{સંબંધાત્મક નિદર્શ}માં જગતોનો ગણ હોય છે, જે દરેક કર્તા
માટેના એક દ્વિસ્થાની સુલભતા સંબંધ વડે જોડાયેલા હોય છે; સાથે
કયા જગતમાં કયાં પ્રતિજ્ઞાપન ચલો સત્ય છે તે નક્કી
કરતી નિમણૂક હોય છે.

\begin{defn}
  બહુ-કર્તા જ્ઞાનસંબંધી ભાષાનું \emph{નિદર્શ} એ ત્રિક
  $\mModel{M} = \tuple{W, R, V}$ છે, જ્યાં
  \begin{enumerate}
  \item $W$ એ ``જગતો''નો અખાલી ગણ છે,
  \item દરેક $a \in G$ માટે ${R}_a$ એ~$W$ પરનો દ્વિસ્થાની
  સુલભતા સંબંધ છે, અને
  \item $V$ એવું વિધેય છે જે દરેક પ્રતિજ્ઞાપન ચલ~$p$ને શક્ય જગતોનો ગણ $V(p)$ આપે છે.
  \end{enumerate}
  $R_a ww'$ હોય ત્યારે કહીએ કે $w'$ એ~$w$માંથી
  \emph{કર્તા aને સુલભ} છે. $w \in V(p)$ હોય ત્યારે કહીએ કે
  $p$ એ~$w$માં \emph{સત્ય} છે.
\end{defn}

અહીં કાર્યરીતિ સામાન્ય મોડલ તર્ક જેવી જ છે; ફક્ત વધુ સુલભતા
સંબંધો ઉમેરાયા છે. આપેલા કર્તાના સુલભતા સંબંધને સામાન્ય રીતે
તેની માહિતી-સ્થિતિ વિશે કંઈક દર્શાવતો માનીએ છીએ. દાખલા તરીકે,
$R_a ww'$નો અર્થ ઘણી વાર એવો લઈએ છીએ કે $w'$ એ~$w$માં~$a$ને
મળેલી માહિતી સાથે સુસંગત છે. બીજા શબ્દોમાં, $w$માં તે
જગત~$w$ અને જગત~$w'$ વચ્ચે ભેદ કરી શકતો નથી.
% END OLP-0485
\endgroup

\begingroup
\makeatletter\def\ol@sectioncs{section}\makeatother

% BEGIN OLP-0486
\olfileid{aml}{el}{trw}

\olsection{જગતમાં સત્ય}
\phantomsection\label{guunit:OLP-0486}


સામાન્ય મોડલ તર્કની જેમ દરેક જ્ઞાનસંબંધી નિદર્શ એ નક્કી કરે છે કે
તેના કયા જગતમાં કયાં સૂત્રો સત્ય ગણાય. સંબંધાત્મક
અર્થવિજ્ઞાનના મૂળ ખ્યાલ માટે આપણે એ જ સંજ્ઞા વાપરીએ છીએ:
``નિદર્શ~$\mModel{M}$માં સૂત્ર~$!A$ જગત~$w$માં સત્ય છે.''
આ સંબંધ આગમનાત્મક રીતે વ્યાખ્યાયિત થાય છે અને બધા બિનમોડલ
સંક્રિયકો માટે સામાન્ય મોડલ તર્કના કિસ્સા જેવો જ છે.

\begin{defn}\ollabel{defn:mmodels}
  નિદર્શ~$\mModel M$માં \emph{સૂત્ર~$!A$નું~$w$માં સત્ય હોવું},
  સંજ્ઞામાં $\mSat{M}{!A}[w]$, નીચે મુજબ આગમનાત્મક રીતે
  વ્યાખ્યાયિત થાય છે:
  \begin{enumerate}
  \item \indcase{!A}{\lfalse}{$\mSat{M}{\lfalse}[w]$ ક્યારેય નહીં}.
  \item \indcase{!A}{\ltrue}{$\mSat{M}{\ltrue}[w]$ હંમેશા}.
  \item $\mSat{M}{p}[w]$ ત્યારે અને માત્ર ત્યારે જ જ્યારે $w \in V(p)$
  \item \indcase{!A}{\lnot !B}{$\mSat{M}{\indfrm}[w]$ ત્યારે અને માત્ર ત્યારે જ જ્યારે
    $\mSat/{M}{!B}[w]$}.
  \item \indcase{!A}{(!B \land !C)}{$\mSat{M}{\indfrm}[w]$ ત્યારે અને માત્ર ત્યારે જ જ્યારે
    $\mSat{M}{!B}[w]$ અને $\mSat{M}{!C}[w]$}.
  \item \indcase{!A}{(!B \lor !C)}{$\mSat{M}{\indfrm}[w]$ ત્યારે અને માત્ર ત્યારે જ જ્યારે
    $\mSat{M}{!B}[w]$ અથવા $\mSat{M}{!C}[w]$} (અથવા બંને).
  \item \indcase{!A}{(!B \lif !C)}{$\mSat{M}{\indfrm}[w]$ ત્યારે અને માત્ર ત્યારે જ જ્યારે
    $\mSat/{M}{!B}[w]$ અથવા $\mSat{M}{!C}[w]$}.
  \item \indcase{!A}{(!B \liff !C)}{$\mSat{M}{\indfrm}[w]$ ત્યારે અને માત્ર ત્યારે જ જ્યારે
    $\mSat{M}{!B}[w]$ અને $\mSat{M}{!C}[w]$ બંને હોય, અથવા
    $\mSat{M}{!B}[w]$ તથા $\mSat{M}{!C}[w]$ બેમાંથી એકેય ન હોય}.
  \item\ollabel{defn:sub:mmodels-box}
    \indcase{!A}{\Knows_a !B}{$\mSat{M}{\indfrm}[w]$ ત્યારે અને માત્ર ત્યારે જ જ્યારે
    $R_a ww'$ ધરાવતા દરેક $w' \in W$ માટે $\mSat{M}{!B}[w']$}
  \end{enumerate} 
\end{defn}

પરંતુ અહીં આપણા સુલભતા સંબંધો પરનાં નિયંત્રણો વિશે વિચારવું પડે.
કારણ કે \olref{defn:sub:mmodels-box}ની શરત મુજબ, $R_a ww'$
ધરાવતું કોઈ~$w'$ ન હોય ત્યારે~$\Knows_a !B$ એ સૂત્ર
જગત~$w$માં સત્ય છે. આ સામાન્ય મોડલ તર્કની જ શરત છે: કોઈ
જગતનાં અનુગામી જગતો ન હોય, તો ત્યાં બધાં $\Box$-સૂત્રો
રિક્તપણે સત્ય હોય છે. પરંતુ $\Knows$ને \emph{જ્ઞાન} દર્શાવતો
માનીએ, તો આ ઘણું વિસંગત લાગે છે. સામાન્ય રીતે આપણે માનીએ છીએ
કે એવો \emph{કોઈ} સંજોગ નથી જેમાં કર્તા એક જ સમયે $!A$ અને
$\lnot !A$ બંને જાણતો હોય.

એક ઉપાય એ છે કે જ્ઞાનસંબંધી તર્કમાં સુલભતા સંબંધ હંમેશા
\emph{સ્વવાચક} હોય તે સુનિશ્ચિત કરીએ. આ લગભગ એ વિચારને અનુરૂપ
છે કે વાસ્તવિક જગત કર્તાની માહિતી સાથે સુસંગત છે. હકીકતમાં,
જ્ઞાનસંબંધી તર્કપદ્ધતિઓ સામાન્ય રીતે S5 વાપરે છે, પરંતુ
$\Knows_a$ સંબંધ વડે ચોક્કસ શું દર્શાવવું છે તેના આધારે
કેટલીક પદ્ધતિઓ નબળાં તંત્રો વાપરી શકે છે.

\begin{figure}
  \begin{center}
    \begin{tikzpicture}[modal]
      \node[world] (w1) [label={[align=right]right:\mTrue{p}\\ \mFalse{q}}]{$w_1$}; 
      \node[world] (w2) [label={[align=right]right:\mTrue{p}\\ \mTrue{q}},
        above right=of w1]{$w_2$}; 
      \node[world] (w3) [label={[align=right]right:\mFalse{p}\\ \mFalse{q}},
        below right=of w1] {$w_3$};
      \draw[<->] (w1) to node {$a$} (w2);
      \draw[->,loop left] (w1) to node {$a,b$} (w1);
      \draw[<->] (w1) to [swap] node {$b$} (w3);
      \draw[->,loop above] (w2) to node {$a,b$} (w2) ;
      \draw[->,loop below] (w3) to node {$a,b$} (w3) ;
    \end{tikzpicture}
  \end{center}
  \caption{એક સરળ જ્ઞાનસંબંધી નિદર્શ.}
  \ollabel{fig:simple}
\end{figure}

\begin{prob}
  નીચેનામાંથી કયા દાવા
  \olref[aml][el][trw]{fig:simple}માં સત્ય છે તે વિચારો:
  \begin{enumerate}
    \item $\mSat{M}{\lnot q}[w_1]$;
    \item $\mSat{M}{\Knows_a \lnot q}[w_1]$;
    \item $\mSat{M}{\Knows_b \lnot q}[w_1]$;
    \item $\mSat{M}{\Knows_b q \lor \Knows_b \lnot q}[w_2]$;
    \item $\mSat{M}{\Knows_a( \Knows_b q \lor \Knows_b \lnot q)}[w_2]$;
    \item $\mSat{M}{\EKnows_{\{a,b\}} \lnot q}[w_3]$;
    \end{enumerate}
\end{prob}

હવે જગતમાં સત્યની મૂળભૂત વ્યાખ્યા આપી છે, એટલે સામાન્ય મોડલ
તર્કના મોડલ પ્રામાણ્યતા અને તાર્કિક પરિણામ જેવા બીજા
અર્થવિજ્ઞાનલક્ષી ખ્યાલો અહીં પણ સીધા લાગુ પડે છે; ફક્ત મોડલ
સંક્રિયકોનું અર્થઘટન કરવાની આ નવી રીત તેમને લાગુ પડે છે.

હવે સામાન્ય જ્ઞાનના સંક્રિયક~$\CKnows_G$ માટે સત્ય-શરતો પણ આપી
શકીએ છીએ. \olref[sfr][rel][ops]{sec}માંથી યાદ કરો કે સંબંધ~$R$નું
\emph{પરંપરિત સંવરણ}~$R^+$ નીચે મુજબ વ્યાખ્યાયિત થાય છે:
\begin{align*}
  R^+ &= \bigcup_{ n \in \mathbb{N}} R^n, 
\intertext{જ્યાં}
  R^0 & = R \text{ અને}\\ 
  R^{n+1} & = \Setabs{\tuple{x, z}}{\exists y (R^n xy \land Ryz)}.
\end{align*}
પછી $G$ કર્તાઓનો સમૂહ હોય ત્યારે, $R_G = ( \bigcup_{b
\in G} R_b )^+$ને બધા કર્તાઓના સુલભતા સંબંધોના સંઘનું
પરંપરિત સંવરણ તરીકે વ્યાખ્યાયિત કરીએ છીએ.

\begin{defn}
જો $G' \subseteq G$ હોય, તો $\mSat{M}{\CKnows_{G'} !A}[ w]$
ત્યારે અને માત્ર ત્યારે જ જ્યારે $R_{G'} w w'$ ધરાવતા દરેક
$w'$ માટે $\mSat{M}{!A}[w']$.
\end{defn}
% END OLP-0486
\endgroup

\begingroup
\makeatletter\def\ol@sectioncs{section}\makeatother

% BEGIN OLP-0487
\olfileid{aml}{el}{acc}

\olsection{સુલભતા સંબંધો અને જ્ઞાનસંબંધી સિદ્ધાંતો}
\phantomsection\label{guunit:OLP-0487}


સામાન્ય મોડલ તર્કપદ્ધતિઓમાં ફ્રેમ-અનુરૂપતા વિશે આપણે જે જાણીએ
છીએ તેને આધારે, જ્ઞાનસંબંધી અર્થઘટન આપતાં લાક્ષણિક
સૂત્રો કેવાં બને છે તે જોવું યોગ્ય છે. જ્ઞાનસંબંધી
તર્કપદ્ધતિઓનું અર્થઘટન સામાન્ય રીતે S5માં થાય છે એમ આપણે કહી
ચૂક્યા છીએ. હવે જુદા જ્ઞાનસંબંધી સિદ્ધાંતો કેવી રીતે દર્શાવાય
છે અને તેઓ જુદી ફ્રેમ-શરતોને કેવી રીતે અનુરૂપ છે તે જોઈએ.

સામાન્ય મોડલ તર્કમાંથી યાદ કરો કે જુદાં મોડલ સૂત્રો
સુલભતા સંબંધોના જુદા ગુણધર્મોને લક્ષણિત કરતાં હતાં. નીચેના
કોષ્ટકમાં ચોક્કસ જ્ઞાનસંબંધી સિદ્ધાંતોને અનુરૂપ એવા થોડા
ગુણધર્મો પસંદ કર્યા છે.

\begin{table}[t]
    \begin{tabular}{| p{.47\linewidth} || p{.47\linewidth} |}
      \hline
      {\emph{જો $R$ \dots હોય}} & {\emph{તો~$\mModel{M}$માં \dots સત્ય છે:}} \\
      \hline \hline
        
      & $\Knows (p \lif q) \lif (\Knows p \lif \Knows q)$ 
      \hfill \newline{(સંવરણ)} \\
      \hline
      \emph{સ્વવાચક}: $\forall w Rww$  
      & $\Knows p \lif p$ \hfill \newline{(સત્યાનુરૂપતા)} \\
      \hline
      \emph{પરંપરિત}: \newline
      $\forall u \forall v \forall w ((Ruv \land Rvw) \lif Ruw)$ & 
      $\Knows p \lif \Knows \Knows p$ \hfill 
           \newline{(હકારાત્મક આત્મનિરીક્ષણ)} \\
      \hline 
      \emph{યુક્લિડિયન}: \newline
      $\forall w \forall u \forall v ((Rwu \land Rwv) \lif Ruv)$ &  
      $\lnot \Knows p \lif \Knows \neg \Knows p$ \hfill 
        \newline{(નકારાત્મક આત્મનિરીક્ષણ)} \\
      \hline
    \end{tabular}
    \caption{ચાર જ્ઞાનસંબંધી સિદ્ધાંતો.}
    \ollabel{tab:four}
  \end{table} 

$T$ સ્વયંસિદ્ધાંતને અનુરૂપ સત્યાનુરૂપતાને આ સિદ્ધાંતોમાં ઘણી વાર
સૌથી ઓછો વિવાદાસ્પદ ગણવામાં આવે છે, કારણ કે તે કહે છે કે
સૂત્ર જાણીતું હોય તો તે સાચું હોવું જ જોઈએ. સંવરણ તથા
હકારાત્મક અને નકારાત્મક આત્મનિરીક્ષણ ઘણાં વધુ વિવાદાસ્પદ છે.

$K$ સ્વયંસિદ્ધાંતને અનુરૂપ સંવરણનો અર્થ એ છે કે કર્તાનું જ્ઞાન
પ્રેરણ હેઠળ બંધ છે. અમુક કિસ્સાઓમાં આ વિશ્વસનીય લાગે. દાખલા
તરીકે, મને ખબર હોઈ શકે કે હું વિક્ટોરિયામાં હોઉં, તો વૅન્કુવર
ટાપુ પર હોઉં. વિચિત્ર સંશયવાદી પરિસ્થિતિઓને બાજુએ રાખીએ, તો
મને ખબર છે કે હું વિક્ટોરિયામાં છું; તેથી મને વૅન્કુવર ટાપુ
પર હોવાની પણ ખબર છે એવું સૂચવાય છે. આ કિસ્સામાં મારા જ્ઞાનનું
તાર્કિક સંવરણ સહજ લાગે. બીજી તરફ, આપણે આપણા જ્ઞાનમાંથી નીકળતાં
બધાં પરિણામો વિશે હંમેશા વિચારતા નથી; તેથી બીજા કિસ્સાઓમાં
આ સિદ્ધાંત ઓછો સહજ લાગી શકે છે.

હકારાત્મક આત્મનિરીક્ષણને ક્યારેક KK-સિદ્ધાંત કહે છે: હું કંઈક
જાણતો હોઉં, તો મને ખબર હોય કે હું તે જાણું છું. આ 4
સ્વયંસિદ્ધાંતનો જ્ઞાનસંબંધી સમકક્ષ છે. એ જ રીતે નકારાત્મક
આત્મનિરીક્ષણ કહે છે કે હું કંઈક \emph{જાણતો ન હોઉં}, તો મને
ખબર હોય કે હું તે જાણતો નથી; આ 5 સ્વયંસિદ્ધાંતનો સમકક્ષ છે.
\sourcecorrection{OLMD-078}{મૂળમાં બંને સિદ્ધાંતોના વિરોધી
દૃષ્ટાંત માટે વ્યક્તિને હકીકતમાં માહિતી ખબર હોય એવું કહે છે.
નકારાત્મક આત્મનિરીક્ષણને ખંડિત કરવા માહિતી ખરેખર અજાણી હોવી
જરૂરી છે; તેથી નીચે બંને સ્થિતિઓ અલગ કરી છે.}
બંને માટે સામાન્ય વિરોધી દૃષ્ટાંતો મળી શકે: મને કોઈ બાબત ખબર
છે કે નહીં તેની ખાતરી ન હોય. મને તે ખરેખર ખબર હોય તો હકારાત્મક
આત્મનિરીક્ષણ નિષ્ફળ થઈ શકે; અને ખરેખર ખબર ન હોય તો નકારાત્મક
આત્મનિરીક્ષણ નિષ્ફળ થઈ શકે.
% END OLP-0487
\endgroup

\begingroup
\makeatletter\def\ol@sectioncs{section}\makeatother

% BEGIN OLP-0488
\olfileid{aml}{el}{bsd}

\olsection{દ્વિઅનુકરણો}
\phantomsection\label{guunit:OLP-0488}


આપણી જ્ઞાનસંબંધી ભાષાની અભિવ્યક્તિ-શક્તિ વિશેનો એક બાકી પ્રશ્ન
નિદર્શો અને તેમાં સત્ય હોય એવાં સૂત્રો વચ્ચેના સંબંધને લગે છે.
ફ્રેમ-અનુરૂપતાનાં પરિણામોમાં આપણે જોયું છે કે અમુક સૂત્રો
ફ્રેમમાં પ્રામાણ્ય હોય, તો એ ફ્રેમ અમુક ગુણધર્મો ધરાવે છે.
પણ શું આપણી મોડલ ભાષા, દાખલા તરીકે, પોતાને જ જતું સ્વવાચક તીર
ધરાવતા જગત અને એક પછી બીજા જગતમાં જતી અનંત શૃંખલા વચ્ચે ભેદ
કરી શકે? એટલે કે, આ બેમાંથી માત્ર એક જ જગતમાં સત્ય હોય એવું
કોઈ સૂત્ર~$A$ છે?

\emph{દ્વિઅનુકરણ} (બાઇસિમ્યુલેશન) એવો સંબંધ છે જે બે
સંબંધાત્મક નિદર્શો વચ્ચે વ્યાખ્યાયિત કરી શકાય છે અને દર્શાવે
છે કે એમની રચના અસરકારક રીતે સમાન છે. આપણે જોઈશું કે તે
આપણી જ્ઞાનસંબંધી ભાષા વડે ઓળખી શકાય એવી નિદર્શોની એક પ્રકારની
સમતુલ્યતા પકડે છે.

\begin{defn}[દ્વિઅનુકરણ]
$M_1 = \tuple{W_1, R_1, V_1}$ અને $M_2 = \tuple{W_2, R_2, V_2}$
એ બે સંબંધાત્મક નિદર્શો હોય. વળી $\mathcal{R} \subseteq W_1 \times W_2$
દ્વિસ્થાની સંબંધ હોય. દરેક $\tuple{w_1, w_2} \in \mathcal{R}$
માટે નીચેની શરતો પૂરી થાય ત્યારે $\mathcal{R}$ને
\emph{દ્વિઅનુકરણ} કહીએ છીએ:
\sourcecorrection{OLMD-079}{મૂળની બીજી અને ત્રીજી શરતમાં
કર્તાઓનો ગણ $A$ લખાયો છે; આ પ્રકરણમાં તે ગણ $G$ તરીકે
વ્યાખ્યાયિત થયો છે. બંને શરતમાં $G$ મૂક્યો છે.}

\begin{enumerate}
  \item બધા પ્રતિજ્ઞાપન ચલો~$p$ માટે
  $w_1 \in V_1(p)$ ત્યારે અને માત્ર ત્યારે જ જ્યારે $w_2 \in V_2(p)$.
  \item બધા કર્તાઓ $a \in G$ અને જગતો $v_1 \in W_1$ માટે,
  $R_{1_a} w_1 v_1$ હોય તો કોઈ $v_2 \in W_2$ એવું હોય કે
  $R_{2_a} w_2 v_2$ અને $\tuple{v_1, v_2} \in
  \mathcal{R}$.
  \item બધા કર્તાઓ $a \in G$ અને જગતો $v_2 \in W_2$ માટે,
  $R_{2_a} w_2 v_2$ હોય તો કોઈ $v_1 \in W_1$ એવું હોય કે
  $R_{1_a} w_1 v_1$ અને $\tuple{v_1, v_2} \in
  \mathcal{R}$.
\end{enumerate}

$M_1$ અને $M_2$ વચ્ચેના દ્વિઅનુકરણમાં જગતો $w_1$ અને $w_2$
જોડાયેલા હોય, તો $\tuple{M_1, w_1} \leftrightarroweq
\tuple{M_2, w_2}$ પણ લખી શકીએ છીએ અને $\tuple{M_1, w_1}$ તથા
$\tuple{M_2, w_2}$ને \emph{દ્વિઅનુકરણ-સમાન} કહીએ છીએ.
\end{defn}

દ્વિઅનુકરણની જુદી શરતો જુદી બાબતો સુનિશ્ચિત કરે છે.
પહેલી શરત બધા પ્રતિજ્ઞાપન ચલો વિશે સહમતિ
સુનિશ્ચિત કરે છે; તેથી દ્વિઅનુકરણ-સમાન જગતોમાં સમાન
મોડલ-મુક્ત સૂત્રો સત્ય હશે. બીજી અને ત્રીજી શરતોને અનુક્રમે
``આગળ'' અને ``પાછળ''ની શરત કહેવાય છે; તે સુલભતા સંબંધોની
રચના સમાન હોવાની ખાતરી આપે છે.

\begin{thm}
જો $\tuple{M_1, w_1} \leftrightarroweq \tuple{M_2, w_2}$, તો દરેક
સૂત્ર~$!A$ માટે $\mSat{M_1}{!A}[w_1]$ ત્યારે અને માત્ર
ત્યારે જ જ્યારે $\mSat{M_2}{!A}[w_2]$.
\end{thm}

\begin{figure}
  \begin{center}
    \begin{tikzpicture}[modal]
      \node[world] (w1) {$w_1$}; 
      \node[world] (w2) [above left=of w1]{$w_2$}; 
      \node[world] (w3) [above right=of w1] {$w_3$};
      \draw[<->] (w1) to node {$a$} (w2);
      \draw[->,loop below] (w1) to node {$a$} (w1);
      \draw[<->] (w1) to [swap] node {$a$} (w3);
      \draw[->,loop above] (w2) to node {$a$} (w2) ;
      \draw[->,loop above] (w3) to node {$a$} (w3) ;
      
      \node[world](v1)[right of=w1, xshift=1.5in]{$v_1$};
      \node[world](v2)[above of=v1]{$v_2$};
      \draw[<->] (v1) to node {$a$} (v2);
      \draw[->,loop below] (v1) to node {$a$} (v1);
      \draw[->,loop above] (v2) to node {$a$} (v2) ;

      \draw[-,dotted] (w1) to (v1) ;
      \draw[-,dotted,bend left=45] (w2) to (v2) ;
      \draw[-,dotted,bend left=30] (w3) to (v2) ;
    \end{tikzpicture}
  \end{center}
  \caption{બે દ્વિઅનુકરણ-સમાન નિદર્શો.}
  \ollabel{fig:bisimilar}
\end{figure}

\olref{fig:bisimilar}માં દર્શાવેલા બે નિદર્શો એકસરખા નથી,
છતાં જગતો $w_1$ અને~$v_1$ને જોડતું દ્વિઅનુકરણ છે. આ
દ્વિઅનુકરણ $w_2$ અને $w_3$ બંનેને~$v_2$ સાથે પણ જોડે છે:
આપણી મોડલ ભાષામાં વ્યક્ત થઈ શકે એવી કોઈ બાબત ખરેખર એમની
વચ્ચે ભેદ કરી શકતી નથી. જોકે $w_2$ અને~$w_3$માં જુદાં
પ્રતિજ્ઞાપન ચલો સત્ય હોત, તો સ્થિતિ જુદી હોત.
% END OLP-0488
\endgroup

\begingroup
\makeatletter\def\ol@sectioncs{section}\makeatother

% BEGIN OLP-0489
\olfileid{aml}{el}{pal}

\olsection{જાહેર જાહેરાતનો તર્ક}
\phantomsection\label{guunit:OLP-0489}


ગતિશીલ જ્ઞાનસંબંધી તર્કપદ્ધતિઓ સમય જતાં અથવા નવી માહિતી મળતાં
કર્તાઓનું જ્ઞાન કેવી રીતે બદલાય છે તે દર્શાવી શકે છે. તેમાંની
ઘણી પદ્ધતિઓ જ્ઞાનમાં થતા ફેરફારોને માહિતીલક્ષી \emph{ઘટનાઓ}
અથવા \emph{અપડેટો} વડે દર્શાવે છે. સૌથી મૂળભૂત અપડેટ
\emph{જાહેર જાહેરાત} છે: કોઈ સૂત્ર સત્યપૂર્વક જાહેર થાય છે
અને બધા કર્તાઓ એકસાથે એ ઘટનાના સાક્ષી બને છે. આ દર્શાવવા
માટે ભાષા નીચે મુજબ વિસ્તરીએ છીએ.

\begin{defn}
$G$ એ કર્તા-સંકેતોનો ગણ હોય. જાહેર જાહેરાતો ધરાવતી બહુ-કર્તા
જ્ઞાનસંબંધી તર્કની મૂળભૂત ભાષામાં નીચેની બાબતો હોય છે:
\begin{enumerate}
  \item અસત્ય માટેનો વિધાનાત્મક અચલાંક~$\lfalse$.
  \item સત્ય માટેનો વિધાનાત્મક અચલાંક~$\ltrue$.
  \item પ્રતિજ્ઞાપન ચલોનો ગણનીય ગણ: $\Obj
    p_0$, $\Obj p_1$, $\Obj p_2$, \dots
  \item વિધાનાત્મક સંયોજકો: \startycommalist
  \ycomma $\lnot$ (નિષેધ)%
  \ycomma $\land$ (સંયોજન)%
  \ycomma $\lor$ (વિયોજન)%
  \ycomma $\lif$ (પ્રેરણ)%
  \sourcecorrection{OLMD-080}{મૂળના સંયોજકોના કોષ્ટકમાં
  દ્વિશરતી $\liff$ રહી ગયું છે, જ્યારે નીચેની આગમનાત્મક
  વ્યાખ્યામાં તે હાજર છે. અહીં તેને ઉમેર્યું છે.}
  \ycomma $\liff$ (દ્વિમુખી પ્રેરણ).
  \item $a \in G$ હોય ત્યારે જ્ઞાન-સંક્રિયક $\Knows_a$.
  \item $!B$ એ સૂત્ર હોય ત્યારે જાહેર-જાહેરાત
  સંક્રિયક $[!B]$.
\end{enumerate}
\end{defn}

જાહેર-જાહેરાત સંક્રિયક બૉક્સ સંક્રિયકની જેમ કાર્ય કરે છે;
તેથી ભાષાની આગમનાત્મક વ્યાખ્યા નીચે મુજબ છે:

\begin{defn}
  જ્ઞાનસંબંધી ભાષાનાં \emph{સૂત્રો} નીચે મુજબ આગમનાત્મક રીતે
  વ્યાખ્યાયિત થાય છે:
  \begin{enumerate}
  \item $\lfalse$ પરમાણુક સૂત્ર છે.

  \item $\ltrue$ પરમાણુક સૂત્ર છે.

  \item દરેક વિધાનાત્મક ચલ $\Obj p_i$ (પરમાણુક) સૂત્ર છે.

  \item જો $!A$ એ સૂત્ર હોય, તો $\lnot !A$ પણ
  સૂત્ર છે.

  \item જો $!A$ અને $!B$ એ સૂત્રો હોય, તો $(!A \land
  !B)$ પણ સૂત્ર છે.

  \item જો $!A$ અને $!B$ એ સૂત્રો હોય, તો $(!A \lor !B)$
  પણ સૂત્ર છે.

  \item જો $!A$ અને $!B$ એ સૂત્રો હોય, તો $(!A \lif !B)$
  પણ સૂત્ર છે.

  \item જો $!A$ અને $!B$ એ સૂત્રો હોય, તો $(!A \liff !B)$
  પણ સૂત્ર છે.

  \item જો $!A$ એ સૂત્ર હોય અને $a \in G$ હોય, તો $\Knows_a !A$
  પણ સૂત્ર છે.
  
  \item જો $!A$ અને $!B$ એ સૂત્રો હોય, તો $[!A] !B$ પણ
  સૂત્ર છે.

  \item આ સિવાય બીજું કશું સૂત્ર નથી.
\end{enumerate}
\end{defn}

સૂત્ર~$[!A] !B$ને ``$!A$ની સત્ય જાહેરાત થયા પછી $!B$
સત્ય છે'' એમ વાંચવાનું છે. જાહેર જાહેરાતોના સંદર્ભમાં ક્યારેક
સામાન્ય જ્ઞાન વિશે પણ વાત કરવી ઉપયોગી થશે; તેથી ભાષામાં
સામાન્ય જ્ઞાનનો સંક્રિયક~$\CKnows_G !A$ પણ હોઈ શકે છે.
% END OLP-0489
\endgroup

\begingroup
\makeatletter\def\ol@sectioncs{section}\makeatother

% BEGIN OLP-0490
\olfileid{aml}{el}{psm}

\olsection{જાહેર જાહેરાતના તર્કનું અર્થવિજ્ઞાન}
\phantomsection\label{guunit:OLP-0490}


જાહેર જાહેરાતના તર્ક માટેનાં સંબંધાત્મક નિદર્શો જ્ઞાનસંબંધી
તર્કનાં નિદર્શો જેવાં જ છે. પરંતુ જાહેર-જાહેરાત સંક્રિયકનું
અર્થવિજ્ઞાન નવું છે.

\begin{defn}\ollabel{defn:mmodels} નિદર્શ~$\mModel M = \tuple{W, R, V}$માં
  \emph{સૂત્ર~$!A$નું~$w$માં સત્ય હોવું}, સંજ્ઞામાં
  $\mSat{M}{!A}[w]$, નીચે મુજબ આગમનાત્મક રીતે વ્યાખ્યાયિત થાય છે:
  \begin{enumerate}
  \item \indcase{!A}{\lfalse}{$\mSat{M}{\lfalse}[w]$ ક્યારેય નહીં}.
  \item \indcase{!A}{\ltrue}{$\mSat{M}{\ltrue}[w]$ હંમેશા}.
  \item $\mSat{M}{p}[w]$ ત્યારે અને માત્ર ત્યારે જ જ્યારે $w \in V(p)$
  \item \indcase{!A}{\lnot !B}{$\mSat{M}{\indfrm}[w]$ ત્યારે અને માત્ર ત્યારે જ જ્યારે
    $\mSat/{M}{!B}[w]$}.
  \item \indcase{!A}{(!B \land !C)}{$\mSat{M}{\indfrm}[w]$ ત્યારે અને માત્ર ત્યારે જ જ્યારે
    $\mSat{M}{!B}[w]$ અને $\mSat{M}{!C}[w]$}.
  \item \indcase{!A}{(!B \lor !C)}{$\mSat{M}{\indfrm}[w]$ ત્યારે અને માત્ર ત્યારે જ જ્યારે
    $\mSat{M}{!B}[w]$ અથવા $\mSat{M}{!C}[w]$} (અથવા બંને).
  \item \indcase{!A}{(!B \lif !C)}{$\mSat{M}{\indfrm}[w]$ ત્યારે અને માત્ર ત્યારે જ જ્યારે
    $\mSat/{M}{!B}[w]$ અથવા $\mSat{M}{!C}[w]$}.
  \item \indcase{!A}{(!B \liff !C)}{$\mSat{M}{\indfrm}[w]$ ત્યારે અને માત્ર ત્યારે જ જ્યારે
    $\mSat{M}{!B}[w]$ અને $\mSat{M}{!C}[w]$ બંને હોય, અથવા
    $\mSat{M}{!B}[w]$ તથા $\mSat{M}{!C}[w]$ બેમાંથી એકેય ન હોય}.
  \item\ollabel{defn:sub:mmodels-box}
    \indcase{!A}{\Knows_a !B}{$\mSat{M}{\indfrm}[w]$ ત્યારે અને માત્ર ત્યારે જ જ્યારે
    $R_a ww'$ ધરાવતા દરેક $w' \in W$ માટે $\mSat{M}{!B}[w']$}
  \item\ollabel{defn:sub:mmodels-pal}
    \indcase{!A}{[!B] !C}{$\mSat{M}{\indfrm}[w]$ ત્યારે અને માત્ર ત્યારે જ જ્યારે
    $\mSat{M}{!B}[w]$ હોય તો $\mSat{M \mid !B}{!C}[w]$}

  
  અહીં $\mModel M \mid !B = \tuple{W', R', V'}$ નીચે મુજબ
  વ્યાખ્યાયિત થાય છે:
  \begin{enumerate}
    \item $W' = \Setabs{ u \in W}{\mSat{M}{!B}[u]}$. એટલે
      $\mModel M \mid !B$નાં જગતો, $\mModel M$માં $!B$ સત્ય હોય
      તે જ જગતો છે.
    \item $R'_a = R_a \cap (W' \times W')$. દરેક કર્તાનો સુલભતા
      સંબંધ ફક્ત~$W'$માં બચેલાં જગતો સુધી મર્યાદિત થાય છે.
    \item $V'(p) = \Setabs{ u \in W'}{u \in V(p) }$. એવી જ રીતે,
      જગતોમાં વિધાનાત્મક ચલની સત્ય-નિમણૂક એ જ રહે છે: માહિતીલક્ષી
      ઘટનાઓ વિધાનાત્મક ચલોનાં સત્યમૂલ્યો બદલતી નથી.
  \end{enumerate}
  
  \end{enumerate} 
\end{defn}

જાહેર જાહેરાતના તર્કની વિશિષ્ટતા એ છે કે નિદર્શ~$\mModel M$માં
સૂત્રનું સત્ય હોવું ક્યારેક $\mModel M$ સિવાયના બીજા
નિદર્શને જોઈને જ નક્કી કરી શકાય છે.

આપણું અર્થવિજ્ઞાન જાહેરાતના સંક્રિયકને $\Box$~સંક્રિયક ગણે છે.
તેથી કોઈ જગતમાં સૂત્ર~$!A$ સત્યપૂર્વક જાહેર કરી ન શકાય,
તો ત્યાં $[!A]!B$ રિક્તપણે સત્ય હશે; જેમ અંતિમ જગતોમાં બધાં
$\Box$ સૂત્રો સત્ય હોય છે.
\sourcecorrection{OLMD-081}{મૂળમાં અહીં $[!A]B$ છે, જ્યારે
ભાષાની વ્યાખ્યામાં પરિણામનું સૂત્ર $!B$ છે. એ સંજ્ઞા
જાળવવા $[!A]!B$ મૂક્યું છે.}

\begin{figure}
  \begin{center}
    \begin{tikzpicture}[modal]
      \node[world] (w1) [label={below left:$p, \lnot q$}]{$w_1$}; 
      \node[world] (w2) [left=of w1, label={left:$\lnot p, \lnot q$}]{$w_2$}; 
      \node[world] (w3) [above=of w1, label={right:$p, q$}] {$w_3$};
      \draw[<->] (w1) to [swap] node {$b$} (w2);
      \draw[->,loop below] (w1) to node {$a, b$} (w1);
      \draw[<->] (w1) to [swap] node {$a$} (w3);
      \draw[->,loop above] (w2) to node {$a, b$} (w2) ;
      \draw[->,loop above] (w3) to node {$a, b$} (w3) ;
      
      \node[world](v1)[right of=w1, xshift=1.25in, label={right:$p, \lnot q$}]{$w'_1$};
      \node[world](v3)[above of=v1,label={right:$p,q$}]{$w'_3$};
      \draw[<->] (v1) to node {$a$} (v3);
      \draw[->,loop below] (v1) to node {$a,b$} (v1);
      \draw[->,loop above] (v3) to node {$a,b$} (v3) ;
      
      \node(m)[below=of w1]{$\mModel M$};
      \node(m')[below=of v1]{$\mModel M \mid p$}; 

      \draw[-,dotted] (w1) to node {\footnotesize $p$ની જાહેરાત}(v1) ;
     
    \end{tikzpicture}
  \end{center}
  \caption{$p$ની જાહેર જાહેરાત પહેલાં અને પછી.}
  \ollabel{fig:announcement-example}
\end{figure}

સૂત્રની જાહેર જાહેરાતને નિદર્શનું સંકોચન, એટલે કે
સૂત્ર સત્ય હોય તે જ જગતો સુધી તેને મર્યાદિત કરવું,
એમ સમજી શકીએ છીએ. \olref{fig:announcement-example}માં~$p$ની
જાહેર જાહેરાતની અસરનું ઉદાહરણ છે. તેમાં નોંધવા જેવી વાત એ છે
કે જાહેરાતના પરિણામે કર્તા~$b$ને~$p$ની ખબર પડે છે, જ્યારે
કર્તા~$a$ને નવી ખબર પડતી નથી (કારણ કે $a$ને~$p$ સત્ય હોવાની
પહેલેથી જ ખબર હતી).

વધુ ઔપચારિક રીતે, $\mSat{M}{\lnot \Knows_b p}[w_1]$ છે,
પરંતુ $\mSat{M \mid p}{\Knows_b p}[w'_1]$ છે. તેથી
$\mSat{M}{[p] \Knows_b p}[w_1]$. જાહેરાતના પરિણામ વિશે
એથી પણ પ્રબળ દાવો કરી શકાય છે. હકીકતમાં,
$\mSat{M}{[p]\CKnows_{\{a,b\}} p}[w_1]$ પણ છે. બીજા શબ્દોમાં,
$p$ જાહેર થયા પછી તે \emph{સામાન્ય જ્ઞાન} બને છે.

પણ શું સામાન્ય રીતે પણ એમ જ થાય? શું~$!A$ની સત્ય જાહેરાત
\emph{હંમેશા} $!A$ને સામાન્ય જ્ઞાન બનાવે છે? કદાચ આશ્ચર્યની વાત
છે કે જવાબ ના છે. હકીકતમાં, અમુક સૂત્રો સત્યપૂર્વક
જાહેર કરી શકાય, પરંતુ જાહેરાત પછી તે સત્ય રહેતાં નથી. દાખલા
તરીકે, \olref{fig:announcement-example}માં~$w_1$ ખાતે
$p \land \lnot \Knows_b p$ની જાહેરાતની અસર વિચારો.
\sourcecorrection{OLMD-082}{મૂળમાં $\mModel M \mid p$ અને
$\mModel M \mid (p \land \lnot \Knows_b p)$ એક જ નિદર્શ
હોવાનું કહે છે; આ આકૃતિમાં પહેલામાં બે જગતો અને બીજામાં
એક જ જગત બચે છે. છતાં બંનેમાં મૂળ પ્રથમ જગતમાં જાહેરાત
પછી $b$ને $p$ ખબર પડે છે, તેથી મુખ્ય નિષ્કર્ષ જળવાય છે.}
હકીકતમાં, $\mModel M \mid p$ અને
$\mModel M \mid (p \land \lnot \Knows_b p)$ એક જ નિદર્શ
નથી: બીજામાં ફક્ત પ્રથમ જગત બચે છે, જ્યારે પહેલામાં પ્રથમ
અને ત્રીજું બંને બચે છે. જોકે આપણે અગાઉ નોંધ્યું છે તેમ,
$\mSat{M \mid p}{\Knows_b p}[w'_1]$. તેથી
$\mSat{M \mid (p \land \lnot \Knows_b p)}{\lnot
(p \land \lnot \Knows_b p)}[w'_1]$. આમ આ સૂત્ર
જાહેર થયા પછી અસત્ય બને છે.
% END OLP-0490
\endgroup



\OLEndPartHook
% END OLP-0482
\endgroup


\OLEndPartHook
% END OLP-0475
